

\D{\newpage}

\subsection[Interval kepercayaan $t$ satu sampel]
    {Interval kepercayaan $\pmb{t}$ satu sampel}
\label{oneSampleTConfidenceIntervals}

\index{data!lumba-lumba dan merkuri|(}

Mari kita mulai mengenal penerapan distribusi $t$
dalam konteks contoh tentang kandungan merkuri
di otot lumba-lumba.
%Dolphins are at the top of the oceanic food chain, which causes dangerous substances such as mercury to concentrate in their organs and muscles.
Konsentrasi merkuri yang tinggi merupakan masalah penting
baik bagi lumba-lumba
maupun hewan lain, seperti manusia, yang sesekali memakannya.
\captionsetup{width=86mm}

\begin{figure}[h]
\centering
\Figures[Seekor lumba-lumba Risso tampak muncul ke permukaan air. Bagian di depan wajahnya sebagian besar berwarna putih, kemudian tubuhnya bercorak abu-abu dan putih.]{0.8}{rissosDolphin}{rissosDolphin.jpg}  \\
\addvspace{2mm}
\begin{minipage}{\textwidth}
   \caption[rissosDolphinPic]{Seekor lumba-lumba Risso.\vspace{-1mm} \\
   -----------------------------\vspace{-2mm}\\
   {\footnotesize Foto oleh Mike Baird (\oiRedirect{textbook-bairdphotos_com}{www.bairdphotos.com}). \oiRedirect{textbook-CC_BY_2}{lisensi~CC~BY~2.0}.}\vspace{-8mm}}
   \label{rissosDolphin}
\end{minipage}
\stdvspace{}
\end{figure}
\captionsetup{width=\mycaptionwidth}

Kita akan menentukan interval kepercayaan bagi rata-rata kandungan merkuri dalam otot lumba-lumba dengan menggunakan sampel 19 lumba-lumba Risso dari wilayah Taiji di Jepang. Data diringkas dalam Gambar~\ref{summaryStatsOfHgInMuscleOfRissosDolphins}. Nilai minimum dan maksimum yang teramati dapat digunakan untuk menilai apakah terdapat pencilan yang jelas.

\begin{figure}[h]
\centering
\begin{tabular}{ccc cc}
\hline
$n$ & $\bar{x}$ & $s$ & minimum & maksimum \\
19   & 4.4	  & 2.3  & 1.7	       & 9.2 \\
\hline
\end{tabular}
\caption{Ringkasan kandungan merkuri dalam otot
    19 lumba-lumba Risso dari wilayah Taiji.
    Pengukuran dinyatakan dalam mikrogram merkuri per gram basah
    otot ($\mu$g/g basah).}
\label{summaryStatsOfHgInMuscleOfRissosDolphins}
\end{figure}

\begin{examplewrap}
\begin{nexample}{Apakah syarat independensi dan
    normalitas terpenuhi untuk kumpulan~data ini?}
  Observasi-observasi tersebut merupakan sampel acak sederhana,
  sehingga independensi masuk akal.
  Statistik ringkasan dalam
  Gambar~\ref{summaryStatsOfHgInMuscleOfRissosDolphins}
  tidak menunjukkan pencilan yang jelas karena
  semua observasi berada dalam 2.5 simpangan baku
  dari rata-rata.
  Berdasarkan bukti ini, syarat normalitas
  tampaknya masuk akal.
\end{nexample}
\end{examplewrap}

Dalam model normal, kita menggunakan $z^{\star}$ dan galat baku untuk menentukan lebar interval kepercayaan. Kita sedikit menyesuaikan rumus interval kepercayaan ketika menggunakan distribusi $t$:
\begin{align*}
&\text{estimasi titik} \ \pm\  t^{\star}_{df} \times SE
&&\to
&&\bar{x} \ \pm\  t^{\star}_{df} \times \frac{s}{\sqrt{n}}
\end{align*}
%The sample mean is the point estimate of interest.
%The standard error is computed using $SE = s/\sqrt{n}$.

\begin{examplewrap}
\begin{nexample}{Dengan menggunakan statistik ringkasan dalam
    Gambar~\ref{summaryStatsOfHgInMuscleOfRissosDolphins},
    hitung galat baku bagi rata-rata kandungan merkuri
    pada $n = 19$ lumba-lumba tersebut.}
  Kita memasukkan $s$ dan $n$ ke dalam rumus:
  $
  %\begin{align*}
  SE
    = s / \sqrt{n}
    = 2.3 / \sqrt{19}
    = 0.528
  %\end{align*}
  $.
\end{nexample}
\end{examplewrap}

Nilai $t^{\star}_{df}$ merupakan nilai kritis yang kita peroleh berdasarkan
tingkat kepercayaan dan distribusi $t$ dengan $df$ derajat
kebebasan.
Nilai batas itu dicari dengan cara yang sama seperti pada distribusi
normal: kita mencari $t^{\star}_{df}$ sedemikian rupa sehingga
proporsi distribusi $t$ dengan $df$ derajat
kebebasan yang berada dalam jarak $t^{\star}_{df}$
dari 0 sama dengan tingkat kepercayaan yang dikehendaki.

\begin{examplewrap}
\begin{nexample}{Ketika $n = 19$, berapakah banyak
    derajat kebebasan yang sesuai?
    Cari $t^{\star}_{df}$ untuk banyaknya derajat kebebasan ini
    dan tingkat kepercayaan 95\%}
  Banyaknya derajat kebebasan mudah dihitung:
  $df = n - 1 = 18$.
  
  Dengan perangkat lunak statistik, kita mencari nilai kritis yang membuat
  luas ekor atas sama dengan 2.5\%:
  $t^{\star}_{18} = 2.10$.
  Luas di bawah -2.10 juga akan sama dengan 2.5\%.
  Artinya, 95\% distribusi $t$ dengan $df = 18$
  terletak dalam jarak 2.10 satuan dari~0.
\end{nexample}
\end{examplewrap}

%\begin{onebox}{Degrees of freedom for a single sample}
%If the sample has $n$ observations and we are examining a single mean, then we use the $t$-distribution with $df=n-1$ degrees of freedom.
%\end{onebox}

%In our current example, we should use the $t$-distribution
%with $df=19-1=18$ degrees of freedom.
%We can generally identify $t_{18}^{\star}$
%using statistical software.
%Alternatively, we could use the $t$-table in
%Appendix~\ref{tDistributionTable}.
%Generally the value of $t^{\star}_{df}$ is slightly larger
%than what we would get under the normal model with~$z^{\star}$.

\begin{examplewrap}
\begin{nexample}{Hitung dan tafsirkan interval kepercayaan 95\%
    bagi rata-rata kandungan merkuri pada lumba-lumba Risso.}
  Kita dapat menyusun interval kepercayaan sebagai
  \begin{align*}
  \bar{x} \ \pm\  t^{\star}_{18} \times SE
    \quad \to \quad 4.4 \ \pm\  2.10 \times 0.528
    \quad \to \quad (3.29, 5.51)
  \end{align*}
  Kita yakin 95\% bahwa rata-rata kandungan merkuri dalam otot
  lumba-lumba Risso berada antara 3.29 dan 5.51 $\mu$g/gram basah,
  yang dianggap sangat tinggi.
\end{nexample}
\end{examplewrap}

\index{data!lumba-lumba dan merkuri|)}

\begin{onebox}{Menentukan
    interval kepercayaan $\pmb{\MakeLowercase{t}}$
    bagi rata-rata}
  Berdasarkan sampel yang terdiri atas $n$ observasi saling independen
  dan mendekati normal, interval kepercayaan bagi rata-rata
  populasi adalah
  \begin{align*}
  &\text{estimasi titik} \ \pm\  t^{\star}_{df} \times SE
  &&\to
  &&\bar{x} \ \pm\  t^{\star}_{df} \times \frac{s}{\sqrt{n}}
  \end{align*}
  dengan $\bar{x}$ sebagai rata-rata sampel, $t^{\star}_{df}$
  sesuai dengan tingkat kepercayaan dan derajat kebebasan
  $df$, dan $SE$ sebagai galat baku yang diestimasi dari
  sampel.
\end{onebox}
